\chapter{Analysis and Design}
\section*{Introduction}
The design phase transforms the company context into explicit requirements,
responsibilities, workflows and data structures. The following sections
separate what the platform must do from how the Angular, Spring Boot and
MySQL components implement it.

\section{Functional requirements}
The platform is organized around a module lifecycle. An instructor creates a
module, attaches source documents, optionally generates a quiz, reviews it and
publishes the module. A learner can then consult the published material and
submit the associated quiz once. The result remains pending until the
instructor reviews the answers.

\begin{table}[H]\centering
\caption{Functional requirements matrix}\label{tab:requirements}
\begin{tabularx}{\textwidth}{>{\bfseries}p{0.12\textwidth}p{0.28\textwidth}X}
ID & Function & Acceptance criterion\\\toprule
F01 & Authentication & An existing user can access the interface according to the role.\\
F02 & Module & A module has a title, description, category, priority and state.\\
F03 & Material & A file is attached to one module and can be opened or downloaded.\\
F04 & AI generation & A source document produces an editable quiz draft.\\
F05 & Manual quiz & An instructor can create questions and answer options directly.\\
F06 & Publication & Only a reviewed quiz can be made available to learners.\\
F07 & Submission & A learner submits a quiz only once.\\
F08 & Review & The instructor sees the learner, module and answer details.\\
F09 & Approval & Approval can be granted and later removed.\\
F10 & Administration & Users, modules and grouped submissions are visible to admins.\\
\bottomrule
\end{tabularx}
\end{table}

\section{Actors and principal use cases}
\begin{table}[H]\centering
\caption{Actors and responsibilities}\label{tab:actors}
\begin{tabularx}{\textwidth}{>{\bfseries}p{0.20\textwidth}X}
\toprule Actor & Main responsibilities\\\midrule
Learner & Consult published material, submit the associated quiz once and follow the approval state.\\
Instructor & Create and publish modules, author quizzes, review detailed answers and approve or revoke submissions.\\
Administrator & Manage users, modules and quizzes, and consult grouped operational activity.\\
\bottomrule
\end{tabularx}
\end{table}

\subsection{Non-functional requirements}
In addition to functional behavior, the platform must satisfy the following
quality requirements:
\begin{itemize}
\item \textbf{Security:} access is authenticated with JWT and controlled by
  role-based permissions.
\item \textbf{Consistency:} module deletion removes dependent quizzes,
  materials, progress and results in a controlled transaction.
\item \textbf{Traceability:} learner answers, submission dates and instructor
  approval decisions are persisted.
\item \textbf{Maintainability:} controllers, services, repositories and
  integration clients have separate responsibilities.
\item \textbf{Usability:} each role can identify its next action quickly on a
  responsive interface.
\item \textbf{AI continuity:} manual quiz authoring remains available when the
  external AI provider is unavailable.
\item \textbf{Availability:} learning material remains accessible after a quiz
  has been submitted.
\end{itemize}

\begin{table}[H]\centering
\caption{Non-functional requirements and acceptance evidence}
\label{tab:nonfunctional}
\begin{tabularx}{\textwidth}{>{\bfseries}p{0.25\textwidth}X}
\toprule Requirement & Acceptance evidence\\\midrule
Security & JWT authentication and role-aware API access.\\
Consistency & Transactional module and user deletion workflows.\\
Traceability & Stored answer map, date, learner, quiz and approval state.\\
Maintainability & Layered services, repositories and specialized integrations.\\
Usability & Role-specific dashboards and disabled submitted-quiz action.\\
AI continuity & Manual quiz creation and explicit AI error responses.\\
\bottomrule
\end{tabularx}
\end{table}

\section{Use-case model}
\img{assets/approved-diagrams/01-usecase-general.png}{Global use-case diagram}
The use-case model separates responsibilities while keeping the main training
workflow shared. AI generation is an assistance function, not an automatic
certification function: generated content must remain editable and subject to
human review.

\subsection{Module and material scenario}
The instructor creates a module with a title, description, category and
priority, then attaches one or more PDF, DOCX or PPTX files. The files remain
associated with the module and can be opened or downloaded by a learner once
the module is published. Deleting the module removes its dependent materials.

\subsection{Quiz authoring and publication scenario}
The instructor chooses either manual authoring or AI generation. In the AI
case, extracted document text is sent to the provider and the returned draft
is validated before it is stored. In the manual case, the instructor enters
questions, options and correct answers directly. In both cases, publication
is a separate action so that an unreviewed draft is not exposed to learners.

\subsection{Learner submission scenario}
The learner opens the published module, consults its material and selects one
option per question. The answer map is stored with persistent question
identifiers. The first submission creates a result and sets the assessment
progress to 100 percent; a second submission is rejected and the material
remains available.

\subsection{Instructor approval scenario}
The instructor sees a compact submission row containing the learner, module,
quiz and date. A details action expands the question-by-question answers and
correct answers. The instructor can approve the result or remove the approval
after a mistake. Certification is therefore a human decision rather than a
score-only calculation.

\subsection{Textual description of major use cases}
\begin{table}[H]\centering
\caption{Main use-case scenarios}\label{tab:usecase-scenarios}
\begin{tabularx}{\textwidth}{>{\bfseries}p{0.24\textwidth}X}
\toprule Use case & Nominal and alternative scenarios\\\midrule
Publish module & The instructor creates a module, attaches material and publishes it. If required information is missing or ownership is invalid, publication is rejected.\\
Generate quiz & The instructor selects material and requests a draft. If extraction or the external API fails, an explicit error is returned and manual authoring remains available.\\
Submit quiz & The learner answers and submits once. If a result already exists, the second submission is rejected and the material remains accessible.\\
Review submission & The instructor opens a submission and expands the detailed answers. Approval can be granted or revoked after review.\\
\bottomrule
\end{tabularx}
\end{table}

\section{User-interface conception}
The learner dashboard prioritizes published modules, accessible files and quiz
availability. The instructor dashboard prioritizes creation, publication and
answer review. The administrator dashboard prioritizes counts, grouped
submissions and management actions. These role-oriented interface decisions
are implemented later in Chapter 3; the screenshots there are validation
evidence rather than requirements analysis.

\section{System architecture}
\img{assets/approved-diagrams/06-diagramme-composants.png}{Component architecture}
The Angular client communicates with REST controllers. Controllers delegate to
services, services coordinate repositories and domain entities, and the
document extraction and AI clients are isolated behind dedicated components.
This separation prevents UI concerns from leaking into persistence logic.

\section{Technical conception}
The technical conception uses a layered REST architecture. The controller
receives the request, the service applies the business rule, the repository
persists the entity and the database stores the durable state. The external AI
API is an integration boundary and is not part of the core database model.
\img{assets/approved-diagrams/05d-class-services.png}{Entity--repository--service responsibilities}

\section{Design principles and patterns}
The solution follows a layered architecture and the Single Responsibility
Principle. Controllers map HTTP requests, services implement business rules,
repositories perform persistence queries and entities represent the relational
model. This separation makes the deletion and submission rules testable
without coupling them to Angular templates.

Document processing follows a strategy-like organization: the extraction
responsibility is isolated from quiz generation so that PDF, DOCX and PPTX
support can evolve independently. The quiz lifecycle is treated as a state
transition problem, separating draft, published and submitted behavior.
Dependency inversion is applied to the external AI provider through a
dedicated client boundary, allowing tests to replace the remote service with a
stub.

\section{Security design}
JWT authentication is integrated as the application authentication mechanism.
After a successful login, the server issues a signed token containing the
identity and role claims. The Angular client attaches the bearer token to API
requests, while Spring Security validates the signature and protects
role-specific routes. This design keeps authorization at the API boundary
rather than relying only on hidden buttons in the interface.

\img{assets/approved-diagrams/04-sequence-auth.png}{Authentication sequence}

\begin{table}[H]\centering
\caption{Security responsibilities}\label{tab:security}
\begin{tabularx}{\textwidth}{>{\bfseries}p{0.25\textwidth}X}
\toprule Component & Responsibility\\\midrule
Login service & Verifies credentials and issues a signed JWT.\\
Angular interceptor & Adds the bearer token to API requests.\\
Security filter & Validates signature, expiration and claims.\\
Role mapping & Converts claims into learner, instructor or admin authorities.\\
Protected controller & Rejects unauthorized or forbidden operations.\\
Password storage & Stores a one-way password hash, never the clear password.\\
\bottomrule
\end{tabularx}
\end{table}

\section{Data model and integrity}
The main relationships are User--Course, Course--Material, Course--Quiz,
Quiz--Question and Learner--QuizResult. A quiz belongs to its module so that
the module is the lifecycle owner. Deletion is transactional: results,
progress records, quizzes and materials are removed in a controlled order
before the module itself is deleted. User deletion removes learner evidence
and detaches owned resources where required.

\begin{table}[H]\centering
\caption{Persistent entity dictionary}\label{tab:data-dictionary}
\begin{tabularx}{\textwidth}{>{\bfseries}p{0.20\textwidth}p{0.28\textwidth}X}
\toprule Entity & Main fields & Responsibility\\\midrule
User & id, name, email, role & Learner, instructor and administrator identity.\\
Course & id, title, status, instructor & Training module and lifecycle owner.\\
Material & id, file name, course & Uploaded PDF, DOCX or PPTX resource.\\
Quiz & id, title, course, published & Assessment attached to one module.\\
Question & id, text, options, correct answer & Question content and grading key.\\
QuizResult & learner, quiz, answers, approved & Submission evidence and approval state.\\
Learner progress & learner, course, progress & Course completion tracking.\\
\bottomrule
\end{tabularx}
\end{table}

\img{assets/approved-diagrams/05c-class-quiz-validation.png}{Quiz validation and lifecycle classes}

\section{Quiz lifecycle}
\img{assets/approved-diagrams/10-sequence-learner-takes-submits-quiz.png}{Learner takes and submits a quiz}
The lifecycle distinguishes draft, published and submitted states. A learner
can continue to consult learning material after submission, but the quiz
action becomes unavailable. The approval state is separate from the score:
submission means evidence was received, while approval means an instructor has
reviewed the evidence.

\section{AI design}
\img{assets/approved-diagrams/02-usecase-detailed-generate-quiz.png}{Detailed AI quiz-generation use case}
\img{assets/approved-diagrams/03-sequence-generate-quiz.png}{AI quiz-generation sequence}
The sequence diagram follows the requested analysis notation: the user
initiates the action through the Angular frontend, the controller reads the
material from the database, and the controller calls the external API. The
returned draft is validated and stored as an unpublished quiz. Invalid or
empty responses are reported explicitly and never bypass instructor
publication.

\img{assets/approved-diagrams/05a-class-ai.png}{AI client and generation classes}
\img{assets/approved-diagrams/05b-class-extract.png}{Document extraction classes}

\section*{Conclusion}
The analysis establishes a traceable relationship between company needs,
quality requirements, security, persistence and implementation boundaries.
The next chapter describes how these decisions were realized and validated in
the backend, frontend, dashboards and API workflows.
